\section{Openness and undecidability: why life must be interactive prediction}
\label{sec:openness}

Computational teleology emphasizes a structural boundary: in universal computational substrates (including universal QCA-like dynamics), many future properties are undecidable or infeasible to predict within bounded resources \cite{Turing1936,Rice1953,Ma2025CT}.
Therefore, biological prediction cannot be ``full simulation'' of the environment; it must be interactive, approximate, and continuously updated.

\paragraph{Online correction near the undecidability boundary.}
In the AEC language, a system maintains survival only if it sustains the inequality \eqref{eq:survival_basic}.
When environmental novelty exceeds the compressive capacity of the internal model, the predictive information rate $\dot I_{\mathrm{pred}}$ drops.
To restore viability, the agent has limited options:
\begin{itemize}
  \item increase exploration/measurement effort (raising dissipation and possibly increasing $\dot I_{\mathrm{pred}}$),
  \item simplify the task by lowering effective resolution (sacrificing precision to regain predictability),
  \item reconfigure architecture (paying geometric Landauer costs) to support faster model updating.
\end{itemize}
These options are precisely the trade-offs seen in biological adaptation: stress responses increase metabolic costs, sensory systems adjust gain and resolution, and learning reorganizes circuits at energetic expense.
These are standard features of cellular and neural adaptation mechanisms \cite{Alberts2014,Friston2010}.

\paragraph{Intelligence as rapid model reconfiguration.}
Within this view, ``intelligence'' is not a separate ontological entity; it is an AEC capability for rapid interactive model updates that maintain $\eta_{\mathrm{pred}}$ when prediction degrades.
This provides an interface-level explanation of why living systems exhibit exploration, plasticity, and multi-scale correction: such features are forced by readout constraints and by the presence of hard computational limits on prediction.


